\documentclass[11pt]{article}
\usepackage[a4paper,margin=1in]{geometry}
\usepackage{amsmath,amssymb,mathtools,bm}
\usepackage{enumitem,booktabs,array}
\usepackage{tcolorbox}
\usepackage{hyperref}
\usepackage[protrusion=true,expansion=false]{microtype}
\hypersetup{colorlinks=true,linkcolor=blue,urlcolor=blue,citecolor=blue}
\setlist[itemize]{topsep=2pt,itemsep=2pt}
\setlist[enumerate]{topsep=2pt,itemsep=2pt}
\newtcolorbox{keybox}{colback=blue!3!white,colframe=blue!40!black,boxrule=0.4pt,arc=1mm}
\newtcolorbox{alertbox}{colback=red!3!white,colframe=red!40!black,boxrule=0.4pt,arc=1mm}
\newtcolorbox{answerbox}{colback=green!3!white,colframe=green!40!black,boxrule=0.4pt,arc=1mm}
\newtcolorbox{drillbox}{colback=yellow!5!white,colframe=orange!60!black,boxrule=0.4pt,arc=1mm}
\newcommand{\EV}{\mathbb{E}}
\newcommand{\Var}{\mathrm{Var}}
\newcommand{\Cov}{\mathrm{Cov}}
\newcommand{\blank}[1]{\underline{\hspace{#1}}}

\title{W11-03: Chava (2014, MS) \\
\large ``Environmental Externalities and Cost of Capital'' --- Active Recall Workbook}
\author{Banking \& Risk Management --- Exam Prep}
\date{\today}
\begin{document}\maketitle

\begin{keybox}
\textbf{Paper in one sentence.} Firms with environmental concerns (KLD-flagged: climate change, hazardous chemicals, land use, operations, other) have significantly higher cost of equity and cost of debt, fewer institutional owners, and smaller loan syndicates --- linking environmental externalities to both equity and debt-market pricing.

\textbf{Placement.} Extends Sharfman-Fernando (W11-01) from equity-only to debt and from cross-section to panel; a pillar of the ESG-cost-of-capital literature.
\end{keybox}

\section*{Round 1: Complete Walk-Through}

\subsection*{1.1 Environmental concern measures}
Using KLD data, Chava identifies six distinct environmental concerns: climate change, hazardous chemicals, ozone-depleting products, substantial emissions, regulatory problems, other. Panel of US firms 1991--2007.

\subsection*{1.2 Cost-of-capital measures}
\begin{itemize}
\item Cost of equity: implied from analysts' forecasts (Gebhardt-Lee-Swaminathan, Claus-Thomas).
\item Cost of debt: interest spread on new loans (DealScan syndicated-loan data).
\item Holdings/syndicate: institutional ownership, number of lenders in syndicate.
\end{itemize}

\subsection*{1.3 Empirical spec}
Panel regression with firm and year FE:
\[
\text{CostOfCapital}_{it}=\alpha_i+\delta_t+\beta\cdot\text{EnvConcern}_{it}+X_{it}'\gamma+\varepsilon_{it}.
\]

\subsection*{1.4 Headline results}
\begin{itemize}
\item \textbf{Equity.} Firms with any environmental concern have $\sim$0.5--1 pp higher implied cost of equity.
\item \textbf{Debt.} Loan spreads rise by $\sim$10--20 bps with environmental concerns.
\item \textbf{Lenders.} Fewer lenders willing to participate; syndicate size shrinks.
\item \textbf{Institutional ownership.} Lower for concerned firms.
\end{itemize}

\subsection*{1.5 Mechanism}
\begin{itemize}
\item Environmental concerns raise tail risk (regulation, litigation, cleanup cost).
\item Norm-constrained investors (pensions, ESG funds) avoid concerned firms.
\item Lender risk aversion to environmental liabilities.
\end{itemize}

\subsection*{1.6 Caveats}
\begin{alertbox}
\begin{itemize}
\item KLD classification is subjective; measurement error likely.
\item Endogeneity: poorly performing firms may cut environmental investment.
\item External validity: sample pre-2008, before modern climate-aware ESG investing.
\end{itemize}
\end{alertbox}

\section*{Round 2: Level 1 Fill-in}
\begin{enumerate}
\item Data source: \blank{1cm} environmental concerns, \blank{1cm}--\blank{1cm}.
\item Cost of equity model: \blank{2cm}-Lee-Swaminathan.
\item Equity cost differential: $\sim$\blank{1cm}--\blank{1cm} pp.
\item Loan-spread differential: $\sim$\blank{1cm}--\blank{1cm} bps.
\item Institutional ownership: \blank{1cm} for concerned firms.
\end{enumerate}

\section*{Round 3: Level 2 Prompts}
\begin{enumerate}
\item Write the panel spec and state the identifying assumption.
\item Explain how environmental concerns raise both equity and debt cost.
\item Compare Chava to Sharfman-Fernando (2008).
\item Design a test for the post-2015 Paris-Agreement era.
\end{enumerate}

\section*{Round 4: Minimal Scaffolding}
From a blank page: (i) data, (ii) cost measures, (iii) spec, (iv) magnitudes, (v) mechanisms.

\section*{Round 5: Blank Page}
Essay: environmental externalities and the cost of capital --- Chava's evidence across debt and equity.

\vspace{6pt}
\begin{drillbox}
\textbf{Speed Drill.} (1) Concern classification source. (2) Sample years. (3) Equity cost delta. (4) Loan-spread delta. (5) Syndicate-size effect.
\end{drillbox}

\section*{Quick Reference \& Answer Key}
\begin{answerbox}
\begin{itemize}
\item \textbf{Measure.} KLD environmental concerns.
\item \textbf{Equity.} Cost of equity $\uparrow$ $\sim$0.5--1 pp.
\item \textbf{Debt.} Loan spread $\uparrow$ 10--20 bps; syndicates shrink.
\item \textbf{Institutions.} Ownership drops.
\end{itemize}
\end{answerbox}

\begin{keybox}
\textbf{30-second exam pitch.} Chava (2014) delivers the most comprehensive test of environmental externalities and the cost of capital. Using KLD environmental-concern indicators (climate change, hazardous chemicals, ozone-depleting products, regulatory problems, emissions) in a panel of US firms 1991--2007, and implied-cost-of-equity models plus Dealscan loan-spread data, he finds: concerned firms have $\sim$0.5--1 pp higher cost of equity, 10--20 bps higher loan spreads, smaller syndicates of lenders, and lower institutional ownership. Mechanisms include environmental tail-risk (regulation, litigation, cleanup), norm-based exclusion by ESG-sensitive investors, and lender risk aversion. The paper extends Sharfman-Fernando (2008) from cross-section equity-only to panel, debt, equity, and holdings, and is a canonical reference for ESG-cost-of-capital work.
\end{keybox}

\end{document}
